\scp{Innovations in the historical phonology}{08-01}
The historical phonology separates \tsc{Taliabu} (Soboyo,
Taliabu, Kadai) not only from nearby Mangole
and Sanana, but also from the rest of \tsc{Sula-Buru}, \tsc{Ambon-Seram} and \tsc{Seram-Tanimbar-Bomberai}.
And some of what separates \tsc{Taliabu} from \tsc{Sula-Buru}
and \tsc{Ambon-Seram} also connects it with neighbouring Banggai
and other \tsc{Celebic} languages,
an affinity that can be seen by a perusal of the data in
\citet{Mead2003a,Mead2003b}, where Proto-Eastern Celebic is presumed to be “the
language ancestral to Saluan-Banggai”.
Key innovations in the historical phonology that link \tsc{Taliabu}
with \tsc{Eastern Celebic} languages are listed below,
and then discussed.

\begin{itemize}[noitemsep]
	\item	*-ay > \it{e}
	\item	*-aw > \it{o}
	\item	*ə > \it{o}
	\item	*y/*z = y
	\item	*j > \it{y, l}, {\0}
	\item	*awa > \it{oa} (otherwise *w = \it{w})
\end{itemize}

\newpage
Selected sound correspondences in the vowels are summarised in \trf{08-02}.
Innovations that \tsc{Taliabu} languages (partially) share with
\tsc{Eastern Celebic} are shaded.

\begin{table}[p]
	\centering\tcp{Vowel reflexes in \tsc{Taliabu} and nearby languages}{08-02}
	\st{6pt}\begin{tabular}{lllllll}\lsptoprule
PMP	&	\hr*ə	&	*-əj	&	*-aw	&	*-ay	&	\hr*u	&	\hr*awa\\ \midrule
Saluan	&	{\hr\hr}o	{\cdr}	&	{\hr\hr}oo	{\cdr}	&	{\hr\hr}o	{\cdr}	&	{\hr\hr}e	{\cdr}	&	{\hr\hr}u	{\cdr}	&	{\hr\hr}oa	{\cdr}\\
Balantak	&	{\hr\hr}o	{\cdr} &	{\hr\hr}e	{\cdr} &	{\hr\hr}o	{\cdr} &	{\hr\hr}e	{\cdr} &	{\hr\hr}u	{\cdr} &	{\hr\hr}oa	{\cdr}\\
Banggai	&	{\hr\hr}o	{\cdr} &	{\hr\hr}oy	{\cdr} &	{\hr\hr}o	{\cdr} &	{\hr\hr}e	{\cdr} &	{\hr\hr}u	{\cdr} &	{\hr\hr}oa	{\cdr}\\
\noalign{\vskip-1.5ex}
\multicolumn{7}{@{}c@{}}{\dashedline{105mm}} \\
\noalign{\vskip-0.5ex}
Soboyo	&	{\hr\hr}o	{\cdr} &	{\hr\hr}e, o	{\cdr} &	{\hr\hr}o	{\cdr} &	{\hr\hr}e	{\cdr} &	{\hr\hr}u	{\cdr} &	{\hr\hr}oa, awa	{\cdr}\\
Taliabu	&	{\hr\hr}o	{\cdr} &	{\hr\hr}e, o	{\cdr} &	{\hr\hr}o	{\cdr} &	{\hr\hr}e	{\cdr} &	{\hr\hr}u	{\cdr} &	{\hr\hr}oa, awa	{\cdr}\\
Kadai		&	{\hr\hr}o	{\cdr} &	{\hr\hr}o	{\cdr} &	{\hr\hr}o	{\cdr} &	{\hr\hr}e	{\cdr} &	{\hr\hr}u	{\cdr} &	{\hr\hr}oa, awa	{\cdr}\\ \midrule
Mangole	&	{\hr\hr}e	&		&	{\hr\hr}a	&	{\hr\hr}a	&	{\hr\hr}u, o	&	{\hr\hr}au\\
Sanana	&	{\hr\hr}e	&	{\hr\hr}i	&	{\hr\hr}a	&	{\hr\hr}a	&	{\hr\hr}u, o	&	{\hr\hr}au\\
Lisela	&	{\hr\hr}e	&	{\hr\hr}e	&	{\hr\hr}a	&	{\hr\hr}a	&	{\hr\hr}u, o	&	{\hr\hr}awa\\
Buru	&	{\hr\hr}e	&	{\hr\hr}e	&	{\hr\hr}a	&	{\hr\hr}a	&	{\hr\hr}u, o	&	{\hr\hr}awa\\
Ambelau	&	{\hr\hr}e	&	{\hr\hr}e	&	{\hr\hr}a	&	{\hr\hr}a	&	{\hr\hr}u, o	&	{\hr\hr}aβa\\ \midrule
\tsc{Ambon-Seram}	&	**ə	&		&	**-aw	&	{\hr\hr}a	&	{\hr\hr}u, o	&	**awa\\
Kayeli (AS)	&	{\hr\hr}e	&	{\hr\hr}e	&	{\hr\hr}a	&	{\hr\hr}a	&	{\hr\hr}u, o	&	{\hr\hr}awa\\
		\lspbottomrule
	\end{tabular}
\end{table}

Selected sound correspondences in the consonants are shown in \trf{08-03}.
Grey shading indicates possible links between \tsc{Taliabu}
and \tsc{Eastern Celebic}, or innovations that
distinguish \tsc{Taliabu} from its neighbours.
Blue shading indicates possible links between \tsc{Taliabu} and \tsc{Sula-Buru}.
There is very little that indicates possible
links between \tsc{Taliabu} and \tsc{Ambon-Seram}.
(\tsc{tr} = traces in the vowels.)

%\begin{table}
	%\centering\tcp{Consonant correspondences in \tsc{Taliabu} and nearby languages}{08-03}
	%\st{2.5pt}\begin{tabular}{llllllllll}\lsptoprule
%PMP			&*b					&*p		&*k												&*ŋ											&\hr*R			&\hr*d				&\hr*z				&\hr*j					&	\hr*y\\ \midrule
%Saluan	&{\hr}b, w	&{\hr}p	&{\hr}k									&{\hr}ŋ									&\hp{**}{\0}&\hp{**}h			&\hp{**}ʤ			&\hp{**}{\0}		&\hp{**}{\0}\\
%Balantak&{\hr}b, w	&{\hr}p	&{\hr}k									&{\hr}ŋ									&\hp{**}r		&\hp{**}r			&\hp{**}s, d	&\hp{**}y, l		&\hp{**}{\0}\\
%Banggai	&{\hr}b			&{\hr}p	&{\hr}k									&{\hr}ŋ									&\hp{**}l		&\hp{**}l			&\hp{**}l, d	&\hp{**}y, l		&\hp{**}y\\
%Soboyo	&{\hr}f			&{\hr}h	&{\hr}k									&{\hr}ŋ									&\hp{**}h		&\hp{**}h, (d)&\hp{**}y			&\hp{**}y, l		&\hp{**}y\\  \midrule
%Taliabu	&{\hr}f			&{\hr}h	&{\hr}k									&{\hr}ŋ									&\hp{**}h		&\hp{**}h, (d)&\hp{**}y			&\hp{**}y, l		&\hp{**}y\\
%Kadai		&{\hr}f			&{\hr}h	&{\hr}k									&{\hr}ŋ									&\hp{**}h, y&							&\hp{**}y			&\hp{**}y				&\hp{**}y\\
%Mangole	&{\hr}f			&{\hr}p	&{\hr}k ({\0} /{\gap\#})&{\hr}ŋ ({\0} /{\gap\#})&\hp{**}{\0}&\hp{**}l			&\hp{**}y			&\hp{**}{\0}, l	&\hp{**}y\\ \midrule
%Sanana	&{\hr}f			&{\hr}p	&{\hr}k ({\0} /{\gap\#})&{\hr}n ({\0} /{\gap\#})&\hp{**}{\0}&\hp{**}h			&\hp{**}y			&\hp{**}{\0}, l	&\hp{**}y\\
%Lisela	&{\hr}ɸ			&{\hr}p	&{\hr}k ({\0} /{\gap\#})&{\hr}ŋ, (n)						&\hp{**}h		&\hp{**}r			&\hp{**}{\0}	&\hp{**}{\0}, l	&\hp{**}{\0}		\\
%Buru		&{\hr}f			&{\hr}p	&{\hr}k ({\0} /{\gap\#})&{\hr}ŋ ({\0} /{\gap\#})&\hp{**}h		&\hp{**}r			&\hp{**}{\0}	&\hp{**}{\0}, l	&\hp{**}{\0}\\
%Ambelau	&{\hr}b, β, h&{\hr}f&{\hr}k, ʔ, {\0}				&{\hr}n, l							&\hp{**}h		&\hp{**}h, l	&\hp{**}l			&\hp{**}{\0}, l	&\hp{**}\tsc{tr}\\ \midrule
%\tsc{Ambon-Seram}	
				%&{\hr}h (p, b)&{\hr}{\0}, h, f	&{\hr}{\0}, ʔ, {(k)}
																										%&{\hr}n									&	**R				&	**r					&	**r					&	**l, {\0}			&\hp{**}\tsc{tr}\\ 
%Kayeli {(AS)}
				%&{\hr}b			&{\hr}h&{\hr}k, {\0}						&{\hr}n									&\hp{**}l		&\hp{**}l			&\hp{**}l			&\hp{**}l, {\0}	&\hp{**}\tsc{tr}\\
		%\lspbottomrule
	%\end{tabular}
%\end{table}

\begin{table}[p]
	\centering\tcp{Consonant reflexes in \tsc{Taliabu} and nearby languages}{08-03}
	\st{3pt}\begin{tabular}{llllllllll}\lsptoprule
PMP			&*b									&*p		&*k				&*ŋ					&\hr*R										&\hr*d				&\hr*z								&\hr*j						&	\hr*y\\ \midrule
Saluan	&{\hr}b, w					&{\hr}p	&{\hr}k	&{\hr}ŋ			&\hp{**}{\0}{\cbl}	&\hp{**}h	{\cdr}&\hp{**}ʤ{\cbl}&\hp{**}{\0}{\cdr}&\hp{**}{\0}{\cdr}\\
Balantak&{\hr}b, w					&{\hr}p	&{\hr}k	&{\hr}ŋ			&\hp{**}r									&\hp{**}r			&\hp{**}s, d					&\hp{**}y, l{\cdr}&\hp{**}{\0}{\cdr}\\
Banggai	&{\hr}b							&{\hr}p	&{\hr}k	&{\hr}ŋ			&\hp{**}l									&\hp{**}l			&\hp{**}l, d					&\hp{**}y, l{\cdr}&\hp{**}y{\cdr}\\
\dshmidrule{10}
Soboyo	&{\hr}f{\cbl}&{\hr}h	&{\hr}k	&{\hr}ŋ			&\hp{**}h		{\cbl}	&\hp{**}h, (d){\cdr}&\hp{**}y{\cbl}&\hp{**}y, l{\cdr}&\hp{**}y{\cdr}\\
Taliabu	&{\hr}f{\cbl}&{\hr}h	&{\hr}k	&{\hr}ŋ			&\hp{**}h		{\cbl}	&\hp{**}h, (d){\cdr}&\hp{**}y{\cbl}&\hp{**}y, l{\cdr}&\hp{**}y{\cdr}\\
Kadai		&{\hr}f{\cbl}&{\hr}h	&{\hr}k	&{\hr}ŋ			&\hp{**}h, y{\cbl}	&				{\cdr}&\hp{**}y{\cbl}&\hp{**}y		{\cdr}&\hp{**}y{\cdr}\\ \midrule
Mangole	&{\hr}f{\cbl}&{\hr}p	&{\hr}k &{\hr}ŋ 		&\hp{**}{\0}{\cbl}	&\hp{**}l			&\hp{**}y{\cbl}&\hp{**}{\0}, l		&\hp{**}y{\cdr}\\
Sanana	&{\hr}f{\cbl}&{\hr}p	&{\hr}k &{\hr}n 		&\hp{**}{\0}{\cbl}	&\hp{**}h			&\hp{**}y{\cbl}&\hp{**}{\0}, l		&\hp{**}y{\cdr}\\
Lisela	&{\hr}ɸ{\cbl}&{\hr}p	&{\hr}k &{\hr}ŋ, (n)&\hp{**}h		{\cbl}	&\hp{**}r			&\hp{**}{\0}	&\hp{**}{\0}, l	&\hp{**}{\0}		\\
Buru		&{\hr}f{\cbl}&{\hr}p	&{\hr}k &{\hr}ŋ 		&\hp{**}h		{\cbl}	&\hp{**}r			&\hp{**}{\0}	&\hp{**}{\0}, l	&\hp{**}{\0}\\
Ambelau	&{\hr}b, β, h				&{\hr}f	&{\hr}k, ʔ, {\0}&{\hr}n, l&\hp{**}h&\hp{**}h, l&\hp{**}l&\hp{**}{\0}, l&\hp{**}\tsc{tr}\\ \midrule
\tsc{Ambon-}&&&&&&&&&\\ 
\tsc{Seram}&\multirow{-2}{*}{{\hr}h (p, b)}
					&\multirow{-2}{*}{{\hr}{\0}, h, f}
					&\multirow{-2}{*}{{\hr}{\0}, ʔ, (k)}
					&\multirow{-2}{*}{{\hr}n}
					&\multirow{-2}{*}{**R}
					&\multirow{-2}{*}{**r}
					&\multirow{-2}{*}{**r}
					&\multirow{-2}{*}{**l, {\0}}
					&\multirow{-2}{*}{\hp{**}\tsc{tr}}\\ 
Kayeli 		&{\hr}b&{\hr}h&{\hr}k, {\0}&{\hr}n&\hp{**}l&\hp{**}l&\hp{**}l&\hp{**}l, {\0}&\hp{**}\tsc{tr}\\
		\lspbottomrule
	\end{tabular}
\end{table}

In the labials, \tsc{Taliabu} *p > \it{h} distinguishes these
languages from \tsc{Sula-Buru} *p = \it{p},
and \tsc{Eastern Celebic} *p = \it{p}.
Though if there were a common ancestry between \tsc{Taliabu}
and either group, *p > \it{h} could be a secondary development.
*b > \it{f} in both \tsc{Taliabu}
and \tsc{Nuclear Sula-Buru} could reflect a shared history,
or could just as easily reflect parallel development,
since it is such a widespread sound change -- and they are in contact.
However, Ambelau retains the voicing
and the obstruent quality in many words: *b > \it{b,
β, h}, thus *b > **f cannot be posited for Proto-Sula-Buru.

The lexically specific split of *j discussed in \crf{ch:CM} looks
simpler in \trf{08-03} than it actually is,
since the split often does not track the same way in the same
words across all the languages (see \trf{07-12} and \trf{08-09}).
So, for example, it is not the case that *j > \it{y} in the top
five languages always tracks to *j > {\0} in the others.
That means that the apparent shared partial merger in \tsc{Taliabu}
of *j/*z/*y > **y is misleading,
since it does not track the same ways with \tsc{Saluan-Banggai} and with \tsc{Sula-Buru}.


